\documentclass{article}
\usepackage{amsmath,amssymb,amsthm}
\usepackage{algorithm}
\usepackage{algpseudocode}
\usepackage{bm}
\usepackage{enumitem}
\usepackage{hyperref}
\newcommand{\E}{\mathbb{E}}
\newcommand{\R}{\mathbb{R}}
\newcommand{\norm}[1]{\left\lVert #1\right\rVert}
\newcommand{\Qop}{\mathcal{Q}}
\newtheorem{theorem}{Theorem}
\newtheorem{lemma}{Lemma}
\begin{document}

\section*{Quantized Stochastic Gradient Descent (QSGD)}

\section*{Mathematical Formulation}
Let $f:\R^{d}\to\R$ be a differentiable (generally non‑convex) objective and let $\{z_t\}_{t\ge 0}$ be i.i.d.\ data samples.
At iteration $t$ a stochastic gradient $g_t$ is computed
\[
g_t \;=\; \nabla f(w_t;z_t)\, .
\]
Instead of transmitting $g_t$ directly, a (possibly lossy) quantization operator
\[
\Qop:\R^{d}\to\R^{d},\qquad \Qop(g) = \text{quantized version of }g,
\]
is applied.  The update rule of QSGD is
\[
\boxed{\,w_{t+1}=w_t-\eta_t\,\Qop(g_t)\,}\tag{1}
\]
where $\eta_t>0$ is the learning rate.  The algorithm only requires that the quantizer be \emph{unbiased}
\[
\E\big[\Qop(g)\mid g\big]=g,\tag{2}
\]
and that its second moment be bounded:
\[
\E\big[\norm{\Qop(g)}^{2}\mid g\big]\le \norm{g}^{2}+\sigma_{q}^{2}.\tag{3}
\]

\section*{Pseudocode}
\begin{algorithm}[H]
\caption{Quantized Stochastic Gradient Descent (QSGD)}
\begin{algorithmic}[1]
\State \textbf{Input:} initial point $w_0\in\R^{d}$, stepsizes $\{\eta_t\}_{t=0}^{T-1}$
\For{$t=0,1,\dots,T-1$}
    \State Sample data $z_t$ (or minibatch) and compute stochastic gradient $g_t=\nabla f(w_t;z_t)$
    \State Quantize: $\tilde g_t\gets \Qop(g_t)$
    \State Update: $w_{t+1}\gets w_t-\eta_t\tilde g_t$
\EndFor
\State \textbf{Output:} $\{w_t\}_{t=0}^{T}$
\end{algorithmic}
\end{algorithm}

\section*{Dry Run Example}
Consider the simple quadratic $f(w)=(w-3)^{2}$, $w\in\R$, with exact gradients (no stochasticity) and a deterministic 1‑bit quantizer that simply rounds each scalar to the nearest integer:
\[
\Qop(g)=\operatorname{round}(g).
\]
Parameters: $w_0=0$, constant stepsize $\eta=0.1$, run for $T=3$ iterations.

\begin{enumerate}[label=Iteration \arabic*:, leftmargin=*, nosep]
\item $t=0$:
    \[
    g_0=\nabla f(w_0)=2(w_0-3)=-6,\qquad
    \tilde g_0=\Qop(g_0)= -6.
    \]
    Update:
    $w_1 = w_0-\eta\tilde g_0 = 0-0.1(-6)=0.6$.
    Objective:
    $f(w_1)=(0.6-3)^{2}=5.76$.
\item $t=1$:
    \[
    g_1=2(w_1-3)=2(0.6-3)=-4.8,\qquad
    \tilde g_1=\Qop(g_1)= -5.
    \]
    Update:
    $w_2 = 0.6-0.1(-5)=1.1$.
    Objective:
    $f(w_2)=(1.1-3)^{2}=3.61$.
\item $t=2$:
    \[
    g_2=2(w_2-3)=2(1.1-3)=-3.8,\qquad
    \tilde g_2=\Qop(g_2)= -4.
    \]
    Update:
    $w_3 = 1.1-0.1(-4)=1.5$.
    Objective:
    $f(w_3)=(1.5-3)^{2}=2.25$.
\end{enumerate}
The iterates move toward the minimiser $w^{\star}=3$, illustrating the effect of quantized updates.

\newpage
\section*{Assumptions}
\begin{enumerate}[label=\textbf{A\arabic*.}, leftmargin=*, nosep]
\item \textbf{Smoothness.} $f$ is $L$‑smooth:
\[
f(u)\le f(v)+\langle\nabla f(v),u-v\rangle+\frac{L}{2}\norm{u-v}^{2},
\qquad\forall u,v\in\R^{d}. \tag{A1}
\]
\item \textbf{Bounded variance of stochastic gradients.} For all $w$,
\[
\E\!\left[\norm{g_t-\nabla f(w_t)}^{2}\mid w_t\right]\le\sigma^{2}. \tag{A2}
\]
\item \textbf{Unbiased quantization with bounded second moment.} The quantizer $\Qop$ satisfies (2)–(3) with a constant $\sigma_q\ge0$ independent of $t$. \tag{A3}
\item \textbf{Stepsize schedule.} Either a constant stepsize $\eta_t=\eta\le\frac{1}{L}$ or a diminishing schedule $\eta_t=\frac{c}{\sqrt{t+1}}$ with $c>0$.
\end{enumerate}

\section*{Main Convergence Theorem}
\begin{theorem}[Convergence of QSGD for non‑convex smooth objectives]\label{thm:main}
Assume \textbf{A1}–\textbf{A4} hold and let $w^{\star}$ be any global minimiser of $f$ (possibly non‑unique).  
Define the average squared gradient norm over $T$ iterations:
\[
\overline{G}_T \;:=\; \frac{1}{T}\sum_{t=0}^{T-1}\E\!\left[\norm{\nabla f(w_t)}^{2}\right].
\]
If a constant stepsize $\eta\le \frac{1}{L}$ is used, then
\[
\overline{G}_T\;\le\;
\frac{2\big(f(w_0)-f(w^{\star})\big)}{\eta T}
\;+\;
L\eta\big(\sigma^{2}+\sigma_{q}^{2}\big).
\tag{4}
\]
Consequently, choosing $\eta = \Theta\!\big(\frac{1}{\sqrt{T}}\big)$ yields
\[
\overline{G}_T = O\!\Big(\frac{1}{\sqrt{T}}\Big).
\]
\end{theorem}

\section*{Theoretical Properties}
\begin{itemize}
\item \textbf{Stability.} The bound (4) is monotone in $\eta$ for $\eta\le 1/L$; too large a stepsize inflates the variance term $L\eta(\sigma^{2}+\sigma_q^{2})$.
\item \textbf{Hyperparameter sensitivity.} The optimal constant stepsize is $\eta^{\star}= \sqrt{\frac{2\big(f(w_0)-f(w^{\star})\big)}{L(\sigma^{2}+\sigma_q^{2})\,T}}$, balancing bias reduction and quantization noise.
\item \textbf{Comparison to unquantized SGD.} Setting $\sigma_q=0$ recovers the standard SGD bound. Quantization only adds the extra variance term $L\eta\sigma_q^{2}$.
\end{itemize}

\section*{Discussion}
The theorem demonstrates that unbiased quantization incurs only an additive variance penalty.  Under the same smoothness and stochastic‑gradient assumptions as classical SGD, QSGD attains the same $O(1/\sqrt{T})$ rate for the expected gradient norm, confirming that communication compression does not deteriorate asymptotic convergence for non‑convex smooth problems.

\newpage
\section*{Preliminaries (Definitions and Lemmas)}
\begin{lemma}[Smoothness inequality]\label{lem:smooth}
If $f$ is $L$‑smooth, then for any $w$ and any vector $u$,
\[
f(w-u)\le f(w)-\langle\nabla f(w),u\rangle+\frac{L}{2}\norm{u}^{2}.
\tag{5}
\]
\end{lemma}
\begin{proof}
Directly from \textbf{A1} with $u\leftarrow -u$.
\end{proof}

\begin{lemma}[Unbiasedness and variance of the quantizer]\label{lem:quant}
For any random vector $g$,
\[
\E\!\left[\Qop(g)\mid g\right]=g,\qquad
\E\!\left[\norm{\Qop(g)}^{2}\mid g\right]\le \norm{g}^{2}+\sigma_{q}^{2}.
\tag{6}
\]
\end{lemma}
\begin{proof}
These are precisely the statements (2)–(3) of assumption \textbf{A3}.
\end{proof}

\section*{Main Proof (Step‑by‑Step Derivation)}
We prove Theorem~\ref{thm:main}.  Throughout the proof, expectations are taken w.r.t.\ all randomness up to iteration $t$ unless otherwise noted.

\bigskip
\noindent\textbf{[Step 1] Apply smoothness to the update (1).}  
Using Lemma~\ref{lem:smooth} with $u=\eta_t\Qop(g_t)$,
\[
f(w_{t+1})\le f(w_t)-\eta_t\langle\nabla f(w_t),\Qop(g_t)\rangle
+\frac{L\eta_t^{2}}{2}\norm{\Qop(g_t)}^{2}. \tag{7}
\]

\bigskip
\noindent\textbf{[Step 2] Take conditional expectation w.r.t.\ $w_t$.}  
Condition on $w_t$ and use Lemma~\ref{lem:quant} together with unbiasedness of $g_t$ (Assumption~\textbf{A2}):
\[
\begin{aligned}
\E\!\left[f(w_{t+1})\mid w_t\right]
&\le f(w_t)-\eta_t\langle\nabla f(w_t),\E[\Qop(g_t)\mid w_t]\rangle
+\frac{L\eta_t^{2}}{2}\E\!\left[\norm{\Qop(g_t)}^{2}\mid w_t\right]\\[4pt]
&= f(w_t)-\eta_t\langle\nabla f(w_t),\E[g_t\mid w_t]\rangle
+\frac{L\eta_t^{2}}{2}\E\!\left[\norm{\Qop(g_t)}^{2}\mid w_t\right] \\[4pt]
&= f(w_t)-\eta_t\norm{\nabla f(w_t)}^{2}
+\frac{L\eta_t^{2}}{2}\E\!\left[\norm{\Qop(g_t)}^{2}\mid w_t\right],
\end{aligned}
\tag{8}
\]
where we used $\E[g_t\mid w_t]=\nabla f(w_t)$.

\bigskip
\noindent\textbf{[Step 3] Bound the second‑moment term.}  
Decompose $g_t$ as $\nabla f(w_t)+(g_t-\nabla f(w_t))$ and use Lemma~\ref{lem:quant}:
\[
\begin{aligned}
\E\!\left[\norm{\Qop(g_t)}^{2}\mid w_t\right]
&\le \E\!\left[\norm{g_t}^{2}\mid w_t\right]+\sigma_q^{2} \qquad\text{(by (6))}\\[4pt]
&= \norm{\nabla f(w_t)}^{2}
   +\E\!\left[\norm{g_t-\nabla f(w_t)}^{2}\mid w_t\right]
   +\sigma_q^{2} \\[4pt]
&\le \norm{\nabla f(w_t)}^{2} + \sigma^{2} + \sigma_q^{2}
\qquad\text{(by \textbf{A2})}.
\end{aligned}
\tag{9}
\]

\bigskip
\noindent\textbf{[Step 4] Plug (9) into (8).}
\[
\E\!\left[f(w_{t+1})\mid w_t\right]
\le f(w_t)-\eta_t\norm{\nabla f(w_t)}^{2}
+\frac{L\eta_t^{2}}{2}\Bigl(\norm{\nabla f(w_t)}^{2}+\sigma^{2}+\sigma_q^{2}\Bigr).
\tag{10}
\]

\bigskip
\noindent\textbf{[Step 5] Rearrange to isolate $\norm{\nabla f(w_t)}^{2}$.}
\[
\begin{aligned}
\E\!\left[f(w_{t+1})\mid w_t\right]
&\le f(w_t)-\Bigl(\eta_t-\frac{L\eta_t^{2}}{2}\Bigr)\norm{\nabla f(w_t)}^{2}
+\frac{L\eta_t^{2}}{2}\bigl(\sigma^{2}+\sigma_q^{2}\bigr).
\end{aligned}
\tag{11}
\]

\bigskip
\noindent\textbf{[Step 6] Choose a constant stepsize $\eta_t\equiv\eta\le 1/L$.}  
Since $\eta\le 1/L$, we have $\eta-\frac{L\eta^{2}}{2}\ge \frac{\eta}{2}$.  Hence
\[
\E\!\left[f(w_{t+1})\mid w_t\right]
\le f(w_t)-\frac{\eta}{2}\norm{\nabla f(w_t)}^{2}
+\frac{L\eta^{2}}{2}\bigl(\sigma^{2}+\sigma_q^{2}\bigr).
\tag{12}
\]

\bigskip
\noindent\textbf{[Step 7] Take full expectation and sum over $t=0,\dots,T-1$.}
\[
\begin{aligned}
\E[f(w_T)] 
&\le f(w_0)-\frac{\eta}{2}\sum_{t=0}^{T-1}\E\!\left[\norm{\nabla f(w_t)}^{2}\right]
+ \frac{L\eta^{2}}{2}T\bigl(\sigma^{2}+\sigma_q^{2}\bigr).
\end{aligned}
\tag{13}
\]

\bigskip
\noindent\textbf{[Step 8] Lower bound $\E[f(w_T)]$ by the global minimum $f(w^{\star})$.}
\[
f(w^{\star})\le \E[f(w_T)].
\tag{14}
\]

\bigskip
\noindent\textbf{[Step 9] Combine (13) and (14) and solve for the average gradient norm.}
\[
\frac{1}{T}\sum_{t=0}^{T-1}\E\!\left[\norm{\nabla f(w_t)}^{2}\right]
\le \frac{2\bigl(f(w_0)-f(w^{\star})\bigr)}{\eta T}
+ L\eta\bigl(\sigma^{2}+\sigma_q^{2}\bigr).
\tag{15}
\]
Equation (15) is exactly the bound (4) claimed in Theorem~\ref{thm:main}.

\bigskip
\noindent\textbf{[Step 10] Optimising the constant stepsize.}  
Set $\eta = \sqrt{\dfrac{2\bigl(f(w_0)-f(w^{\star})\bigr)}{L(\sigma^{2}+\sigma_q^{2})\,T}}$, which satisfies $\eta\le 1/L$ for sufficiently large $T$. Substituting this $\eta$ into (15) yields
\[
\overline{G}_T \le
2\sqrt{\frac{L\bigl(\sigma^{2}+\sigma_q^{2}\bigr)\bigl(f(w_0)-f(w^{\star})\bigr)}{T}}
= O\!\Big(\frac{1}{\sqrt{T}}\Big).
\tag{16}
\]
Thus the claimed $O(1/\sqrt{T})$ rate follows.

\section*{Rate Derivation (With Constants)}
For completeness, write the explicit constant in (4):
\[
\overline{G}_T
\le
\underbrace{\frac{2\Delta_f}{\eta T}}_{\text{bias term}}
\;+\;
\underbrace{L\eta\bigl(\sigma^{2}+\sigma_q^{2}\bigr)}_{\text{variance term}},
\qquad
\Delta_f:=f(w_0)-f(w^{\star}).
\]
Choosing $\eta = \sqrt{\dfrac{2\Delta_f}{L(\sigma^{2}+\sigma_q^{2})\,T}}$ balances the two terms and gives
\[
\overline{G}_T \le 2\sqrt{\frac{L(\sigma^{2}+\sigma_q^{2})\Delta_f}{T}}.
\]

\section*{Special Cases}
\begin{itemize}
\item \textbf{No quantization.} $\sigma_q=0$ reduces (4) to the classical SGD bound for smooth non‑convex functions.
\item \textbf{Deterministic gradients.} $\sigma=0$ (full‑batch) leaves only the quantization variance term $L\eta\sigma_q^{2}$.
\item \textbf{Diminishing stepsize $\eta_t=c/\sqrt{t+1}$.} A similar telescoping argument (omitted for brevity) yields $\overline{G}_T = O(\log T / \sqrt{T})$.
\end{itemize}

\end{document}